\documentclass[a4paper]{article}
\usepackage[a4paper, top = 4cm, bottom = 4cm, left = 3cm, right = 3cm]{geometry}

\usepackage[backend=bibtex, sorting=none]{biblatex}
\addbibresource{../bibliography.bib}

\usepackage{hyperref}
\hypersetup{colorlinks, linkcolor={black}, citecolor={black}, urlcolor={black}}

\usepackage{amsmath}
\usepackage{amsthm}
\usepackage{thmtools}
\usepackage{mathtools}
\usepackage{indentfirst}
\usepackage{amssymb}
\usepackage{pythontex} 
\usepackage{pgf}

\newcommand{\dv}[2]{\frac{d#1}{d#2}}
\newcommand{\dvin}[3]{\frac{d#1}{d#2}\Big\vert_{#3}}
\newcommand{\dvd}[2]{\frac{d^2#1}{d#2^2}}
\newcommand{\dvt}[2]{\frac{d^3#1}{d#2^3}}
\newcommand{\dvf}[2]{\frac{\delta #1}{\delta #2}}
\newcommand{\pdv}[2]{\frac{\partial#1}{\partial#2}}
\newcommand{\pdvd}[3]{\frac{\partial^2 #1}{\partial#2 \partial#3}}
\newcommand{\pdvdu}[2]{\frac{\partial^2 #1}{\partial#2^2}}
\newcommand{\integ}[3]{\int_{#1}^{#2}d#3~}
\newcommand{\poi}[2]{[#1,~#2]}
\newcommand{\poiexp}[2]{\pdv{#1}{q^i} \pdv{#2}{p_i} - \pdv{#2}{q^i} \pdv{#1}{p_i}}

\newcommand{\comm}[2]{[#1,~#2]}
\newcommand{\set}[2]{\{#1\colon#2\}}
\newcommand{\inner}[2]{\langle#1,~#2\rangle}
\newcommand{\av}[1]{\langle#1\rangle}
\newcommand{\avp}[2]{\langle#1\rangle_{#2}}
\newcommand{\ket}[1]{\vert#1\rangle}
\newcommand{\bra}[1]{\langle#1\vert}
\newcommand{\braket}[2]{\langle#1\vert#2\rangle}
\renewcommand\qedsymbol{q.e.d.}
\DeclareMathOperator{\tr}{tr}

\title{On production of gravitational waves: binary stars\\\linespread{1.5}
\small Essay for the General Relativity I exam}
\author{Matteo Zandi}
\date{February 19, 2024}

\begin{document}

\maketitle

\begin{abstract}
    In this essay, we will explicitly derive the quadrupole formula, starting from the linearised Einstein field equations. Furthermore, we will calculate the quadrupole moment in the simple case of two orbiting stars and will use it to find the perturbation of the metric at a point far-away. We will exploit the quadrupole moment of this system to find the power radiated, the variation of the period and the radius over time. Finally, we will consider actual data from the Taylor and Hulse binary system to give an estimate of the amplitude of radiation, power radiated and the variation of the period over time.
\end{abstract}

\section{Introduction}

    It was the year 1916 when Albert Einstein predicted waves solutions, called gravitational waves, for its field equation. It was thought before that a gravitational equivalent of the electromagnetic wave should exist, but a mathematical derivation of that was impossible because there were no gravity equations. Therefore, only after Einstein derived its famous field equations for the general theory of relativity, it was possible to see that they allow indeed a wave solution. However, proving that they were physical it was still in doubt. In fact, Einstein himself, together with one of its collaborator, claimed mistakenly that gravitational waves do not exist in 1936. It was only in 1974 that the first indirect evidence of gravitational waves was discovered: two American physicists Hulse and Taylor studied a binary system composed by a neutron star and a pulsar orbiting around their center of mass, called PSR B1913+16, and found that the experimental observations and theoretical prediction coincide. Needless to say that this discovery allowed them to earn the Nobel Prize in 1993. Since then, technology has been improving so much that in 2015 a direct evidence has been found: by means of a Michelson-Morley type of interferometer, the interference caused by the passage of a gravitational wave has been detected. The system that Hulse and Taylor studied is composed by two neutron stars of equal mass orbiting around their common center of mass with elliptical trajectories. Nevertheless, we will study a simplification of this system, assuming that the orbits are circular. In the next section, we will obtain the mathematical formula necessary to study it: quadrupole formula. In the following sections, we will study in depth the binary stars system by deriving the perturbation of the metric (section 3), the power radiated (section 4) and the variation of the period and radius over time (section 5), so that we can compute actual data in the last section.

\section{Quadrupole formula}

    Our starting point is the linearised Einstein equation in transverse gauge
    \begin{equation}\label{lineq}
        \Box h_{\mu\nu}  (x) = - 16 \pi G_N T_{\mu\nu} (x) ~,
    \end{equation}
    where $\Box = - \partial_t^2 + \nabla^2$ is the d'Alambert operator, $h_{\mu\nu}$ is related to the perturbation of the metric $g_{\mu\nu} = \eta_{\mu\nu} + \overline h_{\mu\nu}$ by 
    \begin{equation*}
        h_{\mu\nu} = \overline h_{\mu\nu} - \frac{1}{2} \eta_{\mu\nu} \overline h
    \end{equation*}
    and $T_{\mu\nu}$ is the energy-momentum tensor. The transverse gauge written in terms of $h_{\mu\nu}$ reads as $\partial^\mu h_{\mu\nu} = 0$. The solution for these equations~\eqref{lineq} can be obtained by means of the Green function $G (x - y)$, which satisfies the Green equation 
    \begin{equation*}
        \Box_x G (x - y) = \delta^4 (x - y) ~.
    \end{equation*}
    The most general solution can be written as a sum of the homogeneous solution $h^{(0)}_{\mu\nu} $, which satisfies the homogeneous equation $\Box h_{\mu\nu}^{(0)}  = 0$, and the inhomogeneous one, written by integrating the Green function:
    \begin{equation*}
        h_{\mu\nu} (x)  = h^{(0)}_{\mu\nu} (x)  - 16 \pi G_N \int d^4 y ~ G (x - y) T_{\mu \nu} (y) ~.
    \end{equation*}
    In fact, if we compute th d'Alambert operator on the solution, we find
    \begin{equation*}
    \begin{aligned}
        \Box_x h_{\mu\nu}  (x) & = \underbrace{\Box h_{\mu\nu (0)} }_0 - 16 \pi G_N \int d^4 y ~ \underbrace{\Box_x G (x - y)}_{\delta^4 (x - y)} T_{\mu \nu} (y) \\ & = - 16 \pi G_N \int d^4 y ~ \delta^4 (x - y) T_{\mu \nu} (y) = - 16 \pi G_N T_{\mu\nu} (x) ~,
    \end{aligned}
    \end{equation*}
    where we have used the properties of the Dirac delta function. Now, let us focus on the inhomogeneous part and neglect the homogeneous one. It is known that the Green function is not unique because it depends on boundary conditions. Since we are interested in the effects of a signal that has propagated in the past, we choose the retarded Green function of the d'Alambertian~\cite{barone}
    \begin{equation*}
        G (x - y) = - \frac{1}{4\pi |\vec x - \vec y|} \theta (x^0 - y^0) \delta (|\vec x - \vec y| - (x^0 - y^0)) ~.
    \end{equation*}
    Therefore, the solution becomes
    \begin{equation*}
    \begin{aligned}
        h_{\mu\nu}  (x) & = - 16 \pi G_N \int d^4 y ~ T_{\mu\nu} (y) G (x - y)  \\ & = - 16 \pi G_N \int d^3 y \int d y^0 ~T_{\mu\nu} (y) \Big (- \frac{1}{4\pi |\vec x - \vec y|} \theta ( x^0 - y^0) \delta (|\vec x - \vec y| - (x^0 - y^0)) \Big ) \\ & = 4 G_N \int d^3 y ~ \frac{T_{\mu\nu} (x^0 - |\vec x - \vec y|, \vec y) }{|\vec x - \vec y|} = 4 G_N \int d^3 y ~ \frac{T_{\mu\nu} (t_r, \vec y) }{|\vec x - \vec y|}~,
    \end{aligned}
    \end{equation*}
    where $t_r = x^0 - |\vec x - \vec y|$ is the retarded time. Physically, it means that the effect of the gravitational disturbance at a point $(x^0, \vec x)$ can be traced back to points on the past light cone $(t_r, \vec x - \vec y)$. The delay in time is due to the finite velocity of gravitational waves. In order to solve this integral, we take the Fourier transform of the perturbation of the metric  
    \begin{equation*}
    \begin{aligned}
        \tilde h_{\mu\nu} (\omega, \vec x) & = \int \frac{dt}{(2\pi)^{1/2}} \exp(i \omega t)  h _{\mu\nu} (t, \vec x) \\ & = 4 G_N \int \frac{dt}{(2\pi)^{1/2}} \exp(i \omega t) \int d^3 y ~ \frac{T_{\mu\nu} (t_r, \vec y)}{|\vec x - \vec y|}  ~,
    \end{aligned}
    \end{equation*}
    which can be developed with a change of variable into $t \rightarrow t_r$ with $t = t_r + |\vec x - \vec y|$ 
    \begin{equation*}
    \begin{aligned}
        \tilde h _{\mu\nu} (\omega, \vec x) & = 4 G_N \int \frac{dt_r}{(2\pi)^{1/2}} d^3 y ~ \exp(i \omega (t_r + |\vec x - \vec y|)) \frac{T_{\mu\nu} (t_r, \vec y)}{|\vec x - \vec y|}  \\ & = 4 G_N \underbrace{\int \frac{dt_r}{(2\pi)^{1/2}} ~ \exp(i \omega t_r) T_{\mu\nu} (t_r, \vec y) }_{\tilde T_{\mu\nu} (\omega, \vec{y})} \int d^3 y ~  \frac{\exp(i \omega |\vec x - \vec y|)}{|\vec x - \vec y|} \\ & = 4 G_N \int d^3 y ~ \tilde T_{\mu\nu} (\omega, \vec{y}) \frac{\exp(i \omega |\vec x - \vec y|)}{|\vec x - \vec y|} ~,
    \end{aligned}
    \end{equation*}
    where we have extracted the Fourier transform of the energy-momentum tensor.

    So far, we have made exact computations. However, we are interested in a simple case, so we consider gravitational waves emitted by a non-relativistic isolated far-away source. We individuate the size of the source to be $r + \delta r$, where $r$ is the (constant) distance of the center of the source from the observer and $\delta r$ is the distance of a part of it from the center. Non-relativistic or slow-moving means that radiation is emitted with low frequency $\delta r \ll 1 / \omega$. Physically, it means that light travels faster than the components of the source so that we can make the following approximation
    \begin{equation*}
        \frac{\exp(i \omega |\vec x - \vec y|)}{|\vec x - \vec y|} \simeq \frac{\exp(i \omega r)}{r} ~, 
    \end{equation*}
    hence, the perturbation of the metric becomes
    \begin{equation}\label{pr1}
        \tilde h_{\mu\nu} (\omega, \vec x) \simeq 4 G_N \frac{\exp(i \omega r)}{r} \int d^3 y ~ \tilde T_{\mu\nu} (\omega, \vec{y}) ~.
    \end{equation}
    Using the transverse gauge $\partial^\mu h_{\mu\nu} (t, \vec x) = 0$, we need only to find the spatial components. In fact, using the inverse Fourier transform, we have 
    \begin{equation*}
    \begin{aligned}
        0 & = \partial^\mu h_{\mu\nu} (t, \vec x) = \partial^\mu \int \frac{d\omega}{(2\pi)^{1/2}} \exp(- i \omega t) \tilde h_{\mu\nu} (\omega, \vec x) \\ & = \partial^0 \int \frac{d\omega}{(2\pi)^{1/2}} \exp(- i \omega t) \tilde h_{0\nu} (\omega, \vec x) + \partial^i \int \frac{d\omega}{(2\pi)^{1/2}} \exp(- i \omega t) \tilde h_{i\nu} (\omega, \vec x) \\ & = i \omega \int \frac{d\omega}{(2\pi)^{1/2}} \exp(- i \omega t) \tilde h_{0\nu} (\omega, \vec x) + \int \frac{d\omega}{(2\pi)^{1/2}} \exp(- i \omega t) \partial^i \tilde h_{i\nu} (\omega, \vec x) ~,
    \end{aligned}
    \end{equation*}
    which implies that 
    \begin{equation*}
        i \omega \tilde h_{0 \nu} + \partial^i \tilde h_{i \nu} = 0 ~, \quad \tilde h_{0\nu} = \frac{i}{\omega} \partial^i \tilde h_{i \nu} ~.
    \end{equation*}
    This means that we only need to find the spatial components and then the temporal components are easily provided by this equation. Henceforth, let us restrict ourselves to the spatial components. Our next goal is to find a usable expression for~\eqref{pr1}
    \begin{equation*}
        \tilde h_{ij} (\omega, \vec x) \simeq 4 G_N \frac{\exp(i \omega r)}{r} \int d^3 y ~ \tilde T_{ij} (\omega, \vec{y}) ~,
    \end{equation*}
    in particular we need to evaluate the integral
    \begin{equation*}
        \int d^3 y ~ \tilde T_{ij} (\omega, \vec{y}) ~.
    \end{equation*}
    In order to do so, we notice the vanishing boundary term
    \begin{equation*}
        0 = \int d^3 y ~ \partial_k (y^i \tilde T^{kj}) = \int d^3 y ~ y^i \partial_k (\tilde T^{kj}) + \int d^3 y ~ \tilde T^{ij}
    \end{equation*}
    and find a useful identity
    \begin{equation}\label{proof5}
        \int d^3 y ~ \tilde T^{ij} = - \int d^3 y ~ y^i \partial_k (\tilde T^{kj}) ~.
    \end{equation}
    Furthermore, we use the continuity equation $\partial^\mu T_{\mu \nu} = 0$ and the inverse Fourier transform, so that 
    \begin{equation*}
    \begin{aligned}
        0 & = \partial^\mu  T_{\mu\nu} (t, \vec x) = \partial^\mu \int \frac{d\omega}{(2\pi)^{1/2}} \exp(- i \omega t) \tilde  T_{\mu\nu} (\omega, \vec x) \\ & = \partial^0 \int \frac{d\omega}{(2\pi)^{1/2}} \exp(- i \omega t) \tilde  T_{0\nu} (\omega, \vec x) + \partial^i \int \frac{d\omega}{(2\pi)^{1/2}} \exp(- i \omega t) \tilde  T_{i\nu} (\omega, \vec x) \\ & = i \omega \int \frac{d\omega}{(2\pi)^{1/2}} \exp(- i \omega t) \tilde  T_{0\nu} (\omega, \vec x) + \int \frac{d\omega}{(2\pi)^{1/2}} \exp(- i \omega t) \partial^i \tilde  T_{i\nu} (\omega, \vec x) ~,
    \end{aligned}
    \end{equation*}
    which implies that 
    \begin{equation}\label{pr3}
        i \omega T^{0 \nu} + \partial_i \tilde T^{i \nu} = 0 ~, \quad \partial_i \tilde T^{i \nu} = i \omega T^{0 \nu} ~.
    \end{equation}
    Therefore, by means of~\eqref{proof5} and~\eqref{pr3}, we find 
    \begin{equation*}
    \begin{aligned}
        \int d^3 y ~ T^{ij} (\omega, \vec y) & = - \int d^3 y ~ y^i \partial_k (\tilde T^{kj}) (\omega, \vec y) = - i \omega \int d^3 y ~ y^i \tilde T^{0j} (\omega, \vec y) \\ & = - \frac{i \omega}{2} \int d^3 y ~ (y^i \tilde T^{0j} (\omega, \vec y) + y^j \tilde T^{0j} (\omega, \vec y)) 
    \end{aligned}
    \end{equation*}
    where we have symmetrised the right-handed side, since the left-hand side $T^{ij}$ is symmetric. Exploiting again~\eqref{proof5} for the first term and~\eqref{pr3} for the second one, we find
    \begin{equation*}
    \begin{aligned}
        \int d^3 y ~ T^{ij} (\omega, \vec y) & = - \frac{i \omega}{2} \int d^3 y ~ \Big (\partial_l (y^i y^l \tilde T^{0j} (\omega, \vec y) ) - y^i y^j \partial_l \tilde T^{0l} (\omega, \vec y) \Big ) \\ & = i^2 \frac{\omega^2}{2} \int d^3 y ~ y^i y^j \tilde T^{00} (\omega, \vec y) = - \frac{\omega^2}{2} \int d^3 y ~ y^i y^j \tilde T^{00} (\omega, \vec y) ~,
    \end{aligned}
    \end{equation*}
    where we have neglected the boundary term. Now, we define the quadrupole moment tensor of the energy density of the source to be
    \begin{equation}\label{quad}
        I^{ij} (t_r) = \int d^3 y ~ y^i y^j T^{00} (t_r, \vec y) 
    \end{equation}
    and evaluate the expression we were looking for
    \begin{equation*}
        \int d^3 y ~ T^{ij} (\omega, \vec y) = - \frac{\omega^2}{2} \int d^3 y ~ y^i y^j \tilde T^{00} = - \frac{\omega^2}{2} \tilde I^{ij} (\omega) ~,
    \end{equation*}
    where $\tilde I_{ij}$ is the Fourier transform of the quadrupole moment. Therefore, the Fourier transform of the perturbation of the metric is
    \begin{equation*}
    \begin{aligned}
        \tilde h _{ij} (\omega, \vec x) & \simeq 4 G_N \frac{\exp(i \omega r)}{r} \int d^3 y ~ \tilde T_{ij} (\omega, \vec{y})  = - \frac{\omega^2}{2} \frac{\exp(i \omega r)}{r} 4 G_N \tilde I^{ij} \\ & = - 2 G_N \omega^2 \frac{\exp(i \omega r)}{r} \tilde I_{ij} (\omega) ~.
    \end{aligned}
    \end{equation*}
    Hence, by the inverse Fourier transform, we obtain
    \begin{equation*}
    \begin{aligned}
        h_{ij} (t, \vec x) & = \int \frac{d\omega}{(2\pi)^{1/2}} \exp(- i \omega t) \tilde h _{ij} (\omega, \vec x) = - 2 G_N \int \frac{d\omega}{(2\pi)^{1/2}} \exp(- i \omega t) \omega^2 \frac{\exp(i \omega r)}{r} \tilde I_{ij} (\omega) \\ & = - \frac{2 G_N}{r} \int \frac{d\omega}{(2\pi)^{1/2}} \omega^2 \exp(- i \omega (t - r)) \tilde I_{ij} (\omega) = - \frac{2 G_N}{r} \int \frac{d\omega}{(2\pi)^{1/2}} \omega^2 \exp(- i \omega t_r) \tilde I_{ij} (\omega) ~.
    \end{aligned}
    \end{equation*}
    Finally, using the Feynman technique
    \begin{equation*}
        \dvd{}{t} I_{ij} (t_r) = \dvd{}{t} \int \frac{d\omega}{(2\pi)^{1/2}} \exp(- i \omega t_r) \tilde I_{ij} (\omega) = - \int \frac{d\omega}{(2\pi)^{1/2}} \omega^2 \exp(- i \omega t_r) \tilde I_{ij} (\omega) ~,
    \end{equation*}
    we derive the so-called quadrupole formula
    \begin{equation}\label{quadform}
        h_{ij} (t, \vec x) = \frac{2 G_N}{r} \dvd{}{t} I_{ij} (t_r) ~.
    \end{equation}
    Its physical meaning is that the produced gravitational wave (perturbation of the metric) is proportional to the second derivative of the quadrupole moment of energy density, measured at a retarded time $t_r$. It was found by Einstein in 1918. Notice that it is different from the electromagnetic case, in which dipole moment is sufficient to produce an electromagnetic wave.

    We conclude this section with a brief comment about the power radiated by the system. It is well known that energy in general relativity is difficult to be defined: we need to find a suitable energy-momentum tensor. By analogy with electromagnetism, the energy-momentum tensor is quadratic in the fields. Therefore, we need to expand at the second order $h_{\mu\nu}$ in order to find what we are looking for. We will not derive here, but the gravitational wave energy-momentum tensor in the transverse traceless gauge is given by
    \begin{equation*}
        t_{\mu\nu} = \frac{1}{32 \pi G_N} \av{\partial_\mu h^{\alpha \beta} \partial_\nu h_{\alpha \beta}} ~.
    \end{equation*}
    Since it is in the transverse traceless gauge, we need to introduce the reduced quadrupole moment 
    \begin{equation}\label{redquad}
        J_{ij} = I_{ij} - \frac{1}{3} \tr I 
    \end{equation}
    and the quadrupole formula becomes 
    \begin{equation*}
        h_{ij}^{TT} (t, \vec x) = \frac{2G_N}{r} \dvd{}{t} J_{ij} (t_r) ~.
    \end{equation*}
    After a series of manipulations~\cite{carroll}, it is possible to prove that the power emitted by the system is given by 
    \begin{equation}\label{pow}
        P = - \frac{G_N}{5} \av{\dvt{}{t} J_{ij} \dvt{}{t} J^{ij}} ~.
    \end{equation}

\section{Binary stars: perturbation of the metric}

    Having found all the mathematical tools, we can apply the formulae we have derived. Consider two point-like stars of equal mass $m$ orbiting around their center of mass $R$ in the plane $xy$. The orbit is circular and can be treated using the Newtonian dynamics. Equating the centrifugal force with the gravitational one
    \begin{equation*}
        \frac{G_N M}{(2 R)^2} = \frac{M v^2}{R} ~, 
    \end{equation*}
    we find the tangential velocity
    \begin{equation*}
        v = \Big ( \frac{G_N M}{4 R} \Big)^{1/2} ~,
    \end{equation*}
    which can be used to calculate the angular frequency of the orbit
    \begin{equation}\label{freq}
        \Omega = \frac{2\pi}{T} = \frac{2 \pi v}{2 \pi R} = \frac{v}{R} = \frac{1}{R} \Big ( \frac{G_N M}{4 R} \Big)^{1/2} = \Big ( \frac{G_N M}{4 R^3} \Big)^{1/2} ~.
    \end{equation}
    Therefore, the explicit path is given by 
    \begin{equation}\label{traj}
        x_a = R \cos (\Omega t) ~, \quad y_a = R \sin (\Omega t) ~, \quad x_b = - R \cos (\Omega t) ~, \quad y_b = - R \sin (\Omega t) ~,
    \end{equation}
    where $a$ and $b$ are label to indicate the two different stars. The energy density of the system is given by 
    \begin{equation*}
        T^{00} (t, \vec x) = M \delta (z) \Big (\delta(x - R \cos (\Omega t))\delta(y - R \sin (\Omega t)) + \delta(x + R \cos (\Omega t))\delta(y + R \sin (\Omega t)) \Big) ~.
    \end{equation*}
    This formula is justified by the fact that we are dealing with point-like objects and all the energy (mass) is located only at a point, mathematically defined by the Dirac delta. The Dirac delta of $z$ constrained the motion on the $xy$-plane. With this component of the energy-momentum tensor, we can explicitly calculate the quadrupole tensor. The components that contain $z$ are zero because of the factor $\delta^3 (z)$. By symmetry, the other components are~\eqref{quad}
    \begin{equation*}
    \begin{aligned}
        I_{11} (t) & = \int d^3 x ~ T^{00} (t, \vec x) x^2 \\ & = \int d^3 x ~ M \delta (z) \Big (\delta(x - R \cos (\Omega t))\delta(y - R \sin (\Omega t)) \\ & \qquad + \delta(x + R \cos (\Omega t))\delta(y + R \sin (\Omega t)) \Big ) x^2 \\ & = M \Big ( \underbrace{\int dx ~ x^2 \delta (x - R \cos (\Omega t))}_{R^2 \cos^2 (\Omega t)} \underbrace{\int dy ~ \delta (y - R \sin (\Omega t))}_1 \underbrace{\int d z ~ \delta (z)}_1 \\ & \qquad + \underbrace{\int dx ~ x^2 \delta (x + R \cos (\Omega t))}_{R^2 \cos^2(\Omega t)}\underbrace{\int dy ~ \delta (y + R \sin (\Omega t))}_1 \underbrace{\int d z ~ \delta (z)}_1 \Big ) \\ & = 2 M R^2 \cos^2 (\Omega t) = M R^2 (1 + \cos (2 \Omega t)) ~,
    \end{aligned}
    \end{equation*}
    \begin{equation*}
    \begin{aligned}
        I_{12} (t) & = I_{21} (t) = \int d^3 x ~ T^{00} (t, \vec x) xy \\ & = \int d^3 x ~ M \delta (z) \Big (\delta(x - R \cos (\Omega t))\delta(y - R \sin (\Omega t)) \\ & \qquad + \delta(x + R \cos (\Omega t))\delta(y + R \sin (\Omega t)) \Big ) x y \\ & = M \Big ( \underbrace{\int dx ~ x \delta (x - R \cos (\Omega t))}_{R \cos (\Omega t)} \underbrace{\int dy ~ y \delta (y - R \sin (\Omega t))}_{R \sin (\Omega t)} \underbrace{\int d z ~ \delta (z)}_1 \\ & \qquad + \underbrace{\int dx ~ x \delta (x + R \cos (\Omega t))}_{- R \cos(\Omega t)}\underbrace{\int dy ~ y \delta (y + R \sin (\Omega t))}_{- R \sin (\Omega t)} \underbrace{\int d z ~ \delta (z)}_1 \Big ) \\ & = 2 M R^2 \cos (\Omega t) \sin (\Omega t) = M R^2 \sin (2 \Omega t) ~,
    \end{aligned}
    \end{equation*}
    \begin{equation*}
    \begin{aligned}
        I_{22} (t) & = \int d^3 x ~ T^{00} (t, \vec x) y^2 \\ & = \int d^3 x ~ M \delta (z) \Big (\delta(x - R \cos (\Omega t))\delta(y - R \sin (\Omega t)) \\ & \qquad + \delta(x + R \cos (\Omega t))\delta(y + R \sin (\Omega t)) \Big ) y^2 \\ & = M \Big ( \underbrace{\int dx ~ \delta (x - R \cos (\Omega t))}_{1} \underbrace{\int dy ~ y^2 \delta (y - R \sin (\Omega t))}_{R^2 \sin^2 (\Omega t)} \underbrace{\int d z ~ \delta (z)}_1 \\ & \qquad + \underbrace{\int dx ~ \delta (x + R \cos (\Omega t))}_{1}\underbrace{\int dy ~ \delta (y + R \sin (\Omega t))}_{R^2 \sin^2(\Omega t)} \underbrace{\int d z ~ \delta (z)}_1 \Big ) \\ & = 2 M R^2 \sin^2 (\Omega t) = M R^2 (1 - \cos (2 \Omega t)) ~.
    \end{aligned}
    \end{equation*}
    Therefore, the quadrupole tensor is 
    \begin{equation}\label{quadsys}
        I_{ij} (t) = MR^2 \begin{bmatrix}
            (1 + \cos (2 \Omega t)) & \sin (2 \Omega t) & 0 \\
            \sin (2 \Omega t) & (1 - \cos (2 \Omega t)) & 0 \\
            0 & 0 & 0 \\
        \end{bmatrix} ~.
    \end{equation}
    Its second derivative is 
    \begin{equation*}
    \begin{aligned}
        \dv{}{t} \dv{}{t} I_{ij} (t) & = MR^2 \dv{}{t} \dv{}{t} \begin{bmatrix}
            (1 + \cos (2 \Omega t)) & \sin (2 \Omega t) & 0 \\
            \sin (2 \Omega t) & (1 - \cos (2 \Omega t)) & 0 \\
            0 & 0 & 0 \\
        \end{bmatrix} \\ & = 2 \Omega MR^2 \dv{}{t} \begin{bmatrix}
            - \sin (2 \Omega t) & \cos (2 \Omega t) & 0 \\
            \cos (2 \Omega t) & \sin (2 \Omega t) & 0 \\
            0 & 0 & 0 \\
        \end{bmatrix} \\ & = 4 \Omega^2 MR^2 \begin{bmatrix}
            - \cos (2 \Omega t) & -\sin (2 \Omega t) & 0 \\
            -\sin (2 \Omega t) & \cos (2 \Omega t) & 0 \\
            0 & 0 & 0 \\
        \end{bmatrix}~.
    \end{aligned}
    \end{equation*}
    Finally, the perturbation of the metric~\eqref{quadform} becomes 
    \begin{equation}\label{met}
    \begin{aligned}
        h_{ij} (t, \vec x) & = \frac{2 G_N}{r} 4 \Omega^2 MR^2 \begin{bmatrix}
            - \cos (2 \Omega t_r) & -\sin (2 \Omega t_r) & 0 \\
            -\sin (2 \Omega t_r) & \cos (2 \Omega t_r) & 0 \\
            0 & 0 & 0 \\
        \end{bmatrix} \\ & = \frac{8 G_N M}{r} \Omega^2 R^2 \begin{bmatrix}
            - \cos (2 \Omega t_r) & -\sin (2 \Omega t_r) & 0 \\
            -\sin (2 \Omega t_r) & \cos (2 \Omega t_r) & 0 \\
            0 & 0 & 0 \\
        \end{bmatrix} ~.
    \end{aligned}
    \end{equation}

\section{Binary stars: power radiated}

    Now, let us calculate the power emitted by our binary system. The trace of the quadrupole moment~\eqref{quadsys} is 
    \begin{equation*}
        \tr I = \tr MR^2 \begin{bmatrix}
            1 + \cos (2 \Omega t) & \sin (2 \Omega t) & 0 \\
            \sin (2 \Omega t) & 1 - \cos (2 \Omega t) & 0 \\
            0 & 0 & 0 \\
        \end{bmatrix} = 2 MR^2 ~.
    \end{equation*}
    Therefore, the reduced quadrupole moment is~\eqref{redquad}
    \begin{equation*}
    \begin{aligned}
        J_{ij} & = I_{ij} - \frac{1}{3} \tr I = MR^2 \begin{bmatrix}
            1 + \cos (2 \Omega t) & \sin (2 \Omega t) & 0 \\
            \sin (2 \Omega t) & 1 - \cos (2 \Omega t) & 0 \\
            0 & 0 & 0 \\
        \end{bmatrix}  - \frac{2}{3} M R^2 \\ & = MR^2 \begin{bmatrix}
            1 + \cos (2 \Omega t) - 2 / 3 & \sin (2 \Omega t) & 0 \\
            \sin (2 \Omega t) & 1 - \cos(2 \Omega t) - 2/3 & 0 \\
            0 & 0 & -2 /3
        \end{bmatrix} \\ & = MR^2 \begin{bmatrix}
            1 / 3 + \cos (2 \Omega t) & \sin (2 \Omega t) & 0 \\
            \sin (2 \Omega t) & 1 / 3 - \cos(2 \Omega t) & 0 \\
            0 & 0 & -2 / 3
        \end{bmatrix} ~.
    \end{aligned}
    \end{equation*}
    Its third derivative is 
    \begin{equation*}
    \begin{aligned}
        \dv{}{t} \dv{}{t} \dv{}{t} J_{ij} (t) & = MR^2 \dv{}{t} \dv{}{t} \dv{}{t}  \begin{bmatrix}
            1 / 3 + \cos (2 \Omega t) & \sin (2 \Omega t) & 0 \\
            \sin (2 \Omega t) & 1 / 3 - \cos(2 \Omega t) & 0 \\
            0 & 0 & -2 / 3
        \end{bmatrix} \\ & = 2 \Omega MR^2 \dv{}{t} \dv{}{t}  \begin{bmatrix}
            - \sin (2 \Omega t) & \cos (2 \Omega t) & 0 \\
            \cos (2 \Omega t) & \sin(2 \Omega t) & 0 \\
            0 & 0 & 0
        \end{bmatrix} \\ & = 4 \Omega^2 MR^2  \dv{}{t}  \begin{bmatrix}
            - \cos (2 \Omega t) & - \sin (2 \Omega t) & 0 \\
            - \sin (2 \Omega t) & \cos(2 \Omega t) & 0 \\
            0 & 0 & 0
        \end{bmatrix} \\ & = 8 \Omega^3 MR^2 \begin{bmatrix}
            \sin (2 \Omega t) & - \cos (2 \Omega t) & 0 \\
            - \cos (2 \Omega t) & - \sin(2 \Omega t) & 0 \\
            0 & 0 & 0
        \end{bmatrix} ~.
    \end{aligned}
    \end{equation*}
    A useful intermediary formula is 
    \begin{equation*}
    \begin{aligned}
        \dddot J_{ij} \dddot J^{ij} & = 64 \Omega^6 M^2 R^4 \Big ( \dddot J_{11}^2 + 2 \dddot J_{12}^2 + \dddot J_{22}^2 \Big) \\ & = 64 \Omega^6 M^2 R^4 \Big ( \sin^2 (2 \Omega t) + 2 \cos^2 (2 \Omega t) + \sin^2 (2 \Omega t) \Big) \\ & = 128 \Omega^6 M^2 R^4 (\cos^2 (2 \Omega t) + \sin^2 (2 \Omega t)) = 128 \Omega^6 M^2 R^4 ~.
    \end{aligned}
    \end{equation*}
    Finally, the power radiated~\eqref{pow} is
    \begin{equation}\label{pow2}
    \begin{aligned}
        P & = - \frac{G_N}{5} \av{\dvt{}{t} J_{ij} \dvt{}{t} J^{ij}} = - \frac{128}{5} G_N \Omega^6 M^2 R^4  \\ & = - \frac{128}{5} G_N \Big (\frac{G_N M}{4 R^3} \Big)^3 M^2 R^4 = - \frac{128}{320} \frac{G_N^4 M^5}{R^5} = - \frac{2}{5} \frac{G_N^4 M^5}{R^5} ~,
    \end{aligned}
    \end{equation}
    where we have used the angular frequency~\eqref{freq}. 

\section{Binary stars: variation of the period}

    The total energy of the system is given by two contributions: the kinetic energy of the two stars and their gravitational potential. The velocity of the stars are the time derivatives of~\eqref{traj}
    \begin{equation*}
        \dot x_a = - \Omega R \sin (\Omega t) ~, \quad y_a = \Omega R \cos (\Omega t) ~, \quad x_b = \Omega R \sin (\Omega t) ~, \quad y_b = - \Omega R \cos (\Omega t) ~,
    \end{equation*}
    Therefore, the total kinetic energy is
    \begin{equation*}
    \begin{aligned}
        T & = T_a + T_b = \frac{M}{2} (\dot x_a^2 + \dot y_a^2) + \frac{M}{2} (\dot x_b^2 + \dot y_b^2) \\ & = \frac{M \Omega^2 R^2}{2} (\underbrace{\sin^2 (\Omega t) + \cos^2(\Omega t)}_1) + \frac{M \Omega^2 R^2}{2} (\underbrace{\sin^2 (\Omega t) + \cos^2(\Omega t)}_1) = M \Omega^2 R^2 ~,
    \end{aligned}
    \end{equation*}
    whereas the potential energy is 
    \begin{equation*}
        V = - \frac{G_N M^2}{2 R} ~. 
    \end{equation*}
    Hence, adding them together, we obtain the total energy
    \begin{equation}\label{ener}
        E = T + V =  M \Omega^2 R^2 - \frac{G_N M^2}{2 R} = M \frac{G_N M}{4 R^3} R^2 - \frac{G_N M^2}{2 R} = - \frac{G_N M^2}{4 R} ~,
    \end{equation}
    where we have used~\eqref{freq}. It is useful to express the total energy in terms of $M$ and $\Omega$, so we invert
    \begin{equation*}
        \Omega^2 = \frac{G_N M}{4 R^3} ~, \quad R^3 = \frac{G_N M}{4 \Omega^2} ~, \quad R = \Big ( \frac{G_N M}{4 \Omega^2} \Big)^{1/3}
    \end{equation*}
    in order to find
    \begin{equation*}
    \begin{aligned}
        E = - \frac{G_N M^2}{4} R = - \frac{G_N M^2}{4} \Big ( \frac{G_N M}{4 \Omega^2} \Big)^{- 1/3} = - 4^{-2/3} G_N^{2/3} M^{5/3} \Omega^{2/3} ~.
    \end{aligned}
    \end{equation*}
    Taking the derivative with respect to $t$
    \begin{equation*}
        \dv{E}{t} = - 4^{-2/3} G_N^{2/3} M^{5/3} \frac{2}{3} \Omega^{-1/3} \dv{\Omega}{t} = \frac{2}{3} \frac{E}{\Omega} \dv{\Omega}{t} 
    \end{equation*}
    and recalling that $P = dE / dt$, we find 
    \begin{equation*}
        P = \frac{2}{3} \frac{E}{\Omega} \dv{\Omega}{t} ~.
    \end{equation*}
    Finally, we compute the period of the orbit using the angular velocity $\Omega = 2 \pi / T$
    \begin{equation*}
        \dv{\Omega}{t} = \dv{}{t} \frac{2\pi}{T} = - \frac{2\pi}{T^2} \dv{T}{t} = - \frac{\Omega}{T} \dv{T}{t} ~,
    \end{equation*}
    hence, we obtain 
    \begin{equation*}
        P = - \frac{2}{3} \frac{E}{\Omega} \frac{\Omega}{T} \dv{T}{t} ~, \quad \dv{T}{t} = - \frac{3 P}{2 E} T ~. 
    \end{equation*}
    Inserting~\eqref{pow2} and~\eqref{ener}, we find
    \begin{equation}\label{var}
        \dv{T}{t} = - \frac{3}{2} \frac{2}{5} \frac{G_N^4 M^5}{R^5} \frac{4 R}{G_N M^2} T = - \frac{12}{5} \frac{G_N^3 M^3}{R^4} T ~,
    \end{equation}
    which is the variation of the period in terms of unit of time.

    The variation of the radius in terms of the period is given by the chain rule
    \begin{equation*}
        \dv{R}{T} = \dv{R}{t} \dv{t}{T} ~.
    \end{equation*}
    Therefore, we need to evaluate the first derivative. Taking the time derivative of the total energy~\eqref{ener}
    \begin{equation*}
        \dv{E}{t} = \frac{G_N M^2}{R^2} \dv{R}{t} = \frac{E}{R} \dv{R}{t} ~,
    \end{equation*}
    and recalling that $P = dE / dt$, we find 
    \begin{equation*}
        P = \frac{E}{R} \dv{R}{t} ~, \quad \dv{R}{t} = \frac{P R}{E} ~.
    \end{equation*}
    Finally, inserting~\eqref{pow2} and~\eqref{ener}, we obtain
    \begin{equation}\label{ret}
        \dv{R}{t} = R \frac{2}{5} \frac{G_N^4 M^5}{R^5} \frac{4 R}{G_N M^2} = \frac{8}{5} \frac{G_N^3 M^3}{R^3} ~.
    \end{equation}
    Putting together~\eqref{var} and~\eqref{ret}, we have
    \begin{equation*}
        \dv{R}{T} = \dv{R}{t} \dv{t}{T} = - \frac{8}{5} \frac{G_N^3 M^3}{R^3} \frac{5}{12} \frac{R^4}{G_N^3 M^3} \frac{1}{T} = - \frac{2}{3} \frac{R}{T} ~.
    \end{equation*}

\section{Hulse and Taylor data}

    In our last section, we will try to estimate orders of magnitude of the physical quantities we have used for the aforementioned Hulse and Taylor system~\cite{schutz}. It is composed by two neutron stars of mass $M = 1.4 ~ M_{\odot}$ orbiting at a period of $T = 7 ~ h ~ 45 ~ min ~ 7 ~ s$. It is distant to us $r = 8 ~kpc$. Converting all data in units of the International System, we find 
    \begin{enumerate}
        \item $T = 27907 ~ s$, 
        \item $\Omega = 2.25 \times 10^{-4} ~ s^{-1}$,
        \item $M = 2.78\times 10^{30} ~ Kg$,
        \item $r = 2.4 \times 10^{20} ~ m$.
    \end{enumerate}
    With these values, we can estimate the amplitude of radiation, the power radiated and the variation of the period for the Hulse and Taylor system. In fact, using~\eqref{met} and reinserting $c$, the amplitude of radiation is of order 
    \begin{equation*}
        A \sim \frac{G_N M \Omega^2 R^2}{c^4 r} \sim \frac{G_N^{5/3} M^{5/3} \Omega^{2/3}}{c^4 r} \sim 10^{-23} ~.
    \end{equation*}
    Using~\eqref{pow2} and reinserting $c$, the power radiated in Watt is
    \begin{equation*}
        P \sim - \frac{G_N^4 M^5}{c^5 R^5} \sim - \frac{G_N^{7/3}M^{10/3}\Omega^{10/3}}{c^5} \sim - 10^{23} ~ W ~.
    \end{equation*}
    Using~\eqref{var} and reinserting $c$, the variation of the period in second per year is
    \begin{equation*}
        \dv{T}{t} \sim - \frac{G_N^3 M^3}{R^4 c^5} T \sim - \frac{G_N^{5/3} M^{5/3} \Omega^{8/3}}{c^5}T \sim - 10^{-14} \sim - 10^{-7} ~ s ~ yr^{-1} ~.
    \end{equation*}
    These results are only estimates for our simple circular model. More correct values can be found by considering that orbits are not circular but ellipses around the center of mass of the stars. In fact, the correct observed value of the variation of the period is of order $10^{-12}$. 

\printbibliography

\end{document}